\section{Conclusions}\label{sec:Conclusion}

In this work, we applied a risk-aware Bayesian optimization campaign (BRAVE) to the Al--V--Cr--Mn--Fe--Co--Ni--Cu FCC HEA system, synthesizing 48 alloys over three closed-loop iterations. We optimized five objectives simultaneously (YS, UTS/YS, strain at UTS, $H_\text{dyn}/H_\text{qs}$, DoP) while exploring 0.12\% of the 27,240-composition feasible space.

\begin{enumerate}

\item \textbf{The highest-performing compositions reside on the FCC phase stability boundary.} Vanadium drives both yield strength ($r = 0.84$, $n = 30$ feasible alloys) and sigma-phase formation ($r = 0.54$, $P = 9 \times 10^{-5}$, $n = 48$). At V = 24~at.\%, the three strongest alloys and three sigma failures share the same compositional point---the definition of a boundary-adjacent optimum. This outcome corresponds to the KKT prediction that constrained optima lie on active boundaries when objective and constraint gradients align~\cite{nocedal2006numerical, ray2009infeasibility}. Although such optima have been observed implicitly in metastability-engineered HEAs~\cite{li2016metastable} and proposed algorithmically~\cite{tian2024boundary, maguire2025good}, they have not, to our knowledge, been demonstrated with statistical significance in a closed-loop experimental campaign.

\item \textbf{BRAVE navigated the boundary where simpler strategies cannot.} Hard filtering would have excluded the V-rich region after the first sigma failures; BBC04 and BBC02 (the two highest-UTS alloys) would never have been synthesized. Conversely, feasibility-blind BO would have repeatedly queried infeasible compositions without learning to avoid them. BRAVE embedded a learned feasibility probability into EHVI, concentrating resources in the V-rich, Ni-rich corner ($<$0.4\% of the feasible space, $P \approx 6.5 \times 10^{-6}$ under random sampling) while reducing the sigma rate from 5/16 to 3/16. Processing failures (1 $\rightarrow$ 0 $\rightarrow$ 6) lie outside the framework's scope; however, despite 9/16 total infeasible alloys in BBC, the 7 feasible compositions included the strongest in the campaign. A computational benchmark (Appendix) confirms that risk-weighted EHVI matches oracle performance with ${\sim}5\times$ fewer wasted evaluations.

\item \textbf{Two compositional regimes emerged from the 8-element space.} A high-strength regime (V = 20--24~at.\%, Cr $\leq$ 4~at.\%, Ni = 36--60~at.\%) produced BBC04 (${\sim}1480$~MPa UTS, 50\% elongation), where Ni stabilizes FCC at extreme V concentrations through VEC effects. A high-ductility regime (V = 8--12~at.\%, Mn = 16~at.\%) achieved UTS/YS up to 4.20 at $>$50\% elongation, consistent with SFE reduction and enhanced strain hardening~\cite{li2016metastable, otto2013influences}. Neither regime was accessible in Campaign~1's 6-element system. SHAP analysis confirmed V as the most influential feature (highest mean $|\text{SHAP}|$), consistent with its strong linear correlation with strength ($r = 0.82$). Ni, ranked fourth by mean $|\text{SHAP}|$, shows near-zero linear correlation with strength ($r = -0.03$) but non-negligible SHAP contributions whose sign reverses with compositional context---consistent with interaction-mediated, non-monotonic influence. Cu correlated negatively with all strength metrics and positively with DoP.

\item \textbf{The campaign avoided the strength--ductility trade-off and compounded gains from Campaign~1.} Median YS increased across iterations while ductility was preserved, and inter-quartile ranges narrowed---indicating convergence. Every single-phase FCC alloy exceeded Campaign~1's best UTS/YS (2.93~\cite{hastings2025accelerated}): minimum 3.07, median 3.79, 9 above 4.0. We attribute this to three compounding factors: (i)~informative priors from Campaign~1 data and the Varvenne--Luque--Curtin model, (ii)~a design space expanded based on Campaign~1's lessons, and (iii)~objectives refined from the finding that hardness alone is an incomplete proxy for mechanical response.

\item \textbf{DoP operated as a simulation-in-the-loop objective.} Calibrated from tensile and SHPB data via a Cowper--Symonds model, DoP tracks YS ($r = -0.94$) but also captures rate-sensitivity effects invisible to quasi-static testing. The strongest alloys (high V) showed lower rate-dependent strengthening, consistent with dislocation-glide-dominated deformation---a trade-off that is invisible without a rate-sensitive objective.

\item \textbf{CALPHAD database fidelity limits the feasibility classifier.} The mid-campaign TCHEA6 $\rightarrow$ TCHEA7 transition reduced but did not eliminate false negatives---the errors that produce failed experiments. The GPC learns the empirical boundary from phase outcomes; however, a poorly calibrated database can exclude high-performing regions or admit too many infeasible compositions before the GPC acquires sufficient training data.

\end{enumerate}

We note several limitations. The campaign optimized composition only; microstructure and processing were not design variables. Variability in grain size (10--35~$\mu$m) and recrystallization response introduced scatter that compositional surrogates cannot capture---several phase-pure alloys showed anomalous tensile behavior traceable to incomplete recrystallization, reducing the signal-to-noise available to the GP models. The DoP model uses a Cowper--Symonds framework without damage criteria, limiting its scope to relative penetration ranking. Furthermore, the GPC (trained on 8 sigma failures) cannot predict the 7 processing failures (i.e., cold-rolling defects) that dominated BBC's budget---a distinct failure mode outside the classifier's design scope.

The three-iteration budget was set by experimental throughput rather than convergence criteria, and the diminishing HV gain (24.2\% in BBB, 4.7\% in BBC) leaves open whether additional iterations would yield substantial further improvement.

In future work, we plan to incorporate microstructural descriptors (e.g., grain size, texture, precipitate fraction) as design variables and extend the feasibility model to processing failures. Testing BRAVE on systems with more complex phase competition (e.g., refractory HEAs with BCC/HCP/sigma boundaries) would assess generalizability. The KKT observation predicts boundary-adjacent optima should recur whenever the performance-driving element also destabilizes the target phase.
